% Solutions to the nine homework problems. Written for this edition; the
% printed notes contain no solutions. Conventions are those of the notes:
% Euclidean signature from Section 5 on, mostly plus metric in Minkowski space.
\appendix
\section{Solutions to the homework problems}\label{sec:solutions}

These solutions use the conventions of the sections where the problems appear.
The link at the end of each solution leads back to the problem.

\tikzset{
  % tree with two proper three-point vertices joined by a full propagator
  pics/gtree/.style={code={%
    \draw[prop] (-0.9,0) -- (0.9,0);
    \draw[prop,fill=white] (0,0) circle (0.2);
    \foreach \a/\p in {135/-0.9,225/-0.9} {\draw[prop] (\p,0) -- ++(\a:1.5) node[vtx] {};
      \draw[prop,fill=white] (\p,0) ++(\a:0.85) circle (0.2);}
    \foreach \a/\p in {45/0.9,-45/0.9} {\draw[prop] (\p,0) -- ++(\a:1.5) node[vtx] {};
      \draw[prop,fill=white] (\p,0) ++(\a:0.85) circle (0.2);}
    \draw[prop,pattern=north east lines] (-0.9,0) circle (0.35);
    \draw[prop,pattern=north east lines] (0.9,0) circle (0.35);}},
}

\subsection{Problem 1}\label{sol:1}

\emph{(a)} Vary $S=\int dt\,L$. In the bulk term, integrate by parts in $t$,
$x_\parallel$ and $z$. The $z$ integration of
$-\partial_z\varphi\,\partial_z\delta\varphi$ leaves the boundary term
$+\partial_z\varphi\,\delta\varphi\big|_{z=0}$, because the fields vanish at
$z\to\infty$. The boundary potential contributes $-V'(\varphi_B)\,\delta\varphi_B$.
So
\[
\delta S=\int dt\int_0^\infty dz\int d^2x_\parallel\,\bigl(-\ddot\varphi+\nabla^2\varphi-m^2\varphi\bigr)\,\delta\varphi
+\int dt\int d^2x_\parallel\,\bigl(\partial_z\varphi\big|_{z=0}-V'(\varphi_B)\bigr)\,\delta\varphi_B\,.
\]
The bulk and boundary variations are independent, so
\[
\ddot\varphi-\nabla^2\varphi+m^2\varphi=0\quad(z>0)\,,\qquad
\partial_z\varphi\big|_{z=0}=V'(\varphi_B)\,.
\]
For $V=0$ this is the Neumann condition.

\emph{(b)} Time translations and translations along $x^1,x^2$ are still
symmetries, so $E$, $P^1$ and $P^2$ are conserved. The boundary breaks
translations in $z$. The energy has to include the boundary term,
\[
E=\int_0^\infty dz\int d^2x_\parallel\,\Bigl[\tfrac12\dot\varphi^2+\tfrac12(\nabla\varphi)^2+\tfrac{m^2}2\varphi^2\Bigr]+\int d^2x_\parallel\,V(\varphi_B)\,.
\]
Using the field equation, the time derivative of the bulk part is
$-\int d^2x_\parallel\,\dot\varphi_B\,\partial_z\varphi_B=-\int d^2x_\parallel\,\dot\varphi_B\,V'(\varphi_B)$,
which cancels $\frac d{dt}\int V(\varphi_B)$. For $P^a=-\int\dot\varphi\,\partial_a\varphi$
($a=1,2$) the same manipulation leaves $-\int d^2x_\parallel\,\partial_aV(\varphi_B)=0$.
For $P^3=-\int\dot\varphi\,\partial_z\varphi$ one finds instead
\[
\frac{dP^3}{dt}=\frac12\int d^2x_\parallel\,\Bigl[\dot\varphi_B^2-(\partial_a\varphi_B)^2-m^2\varphi_B^2+V'(\varphi_B)^2\Bigr]\,,
\]
the force with which the boundary pushes on the field. It does not vanish in
general, even for $V=0$, so $P^3$ is not conserved.

\emph{(c)} With $V=M\varphi_B^2$ the boundary condition is
$\partial_z\varphi=2M\varphi$ at $z=0$. Separate variables,
$\varphi=\Re\bigl[A\,e^{i\kk\cdot x_\parallel-i\omega t}f(z)\bigr]$, so that
$f''=(\kk^2+m^2-\omega^2)f$. A bounded $f$ is either oscillating, $f\propto e^{\pm iqz}$
with $\omega^2=\kk^2+q^2+m^2$, or decaying, $f=e^{-\kappa z}$ with $\kappa>0$. The
decaying one would need $-\kappa=2M$, which is impossible for $M>0$. The oscillating
solutions satisfy the boundary condition in the combination
\[
f_q(z)=q\cos qz+2M\sin qz=\sqrt{q^2+4M^2}\,\cos(qz-\delta_q)\,,\qquad\tan\delta_q=\frac{2M}q\,,
\]
an incident wave plus a reflected wave with a phase shift. The general bounded
solution is the superposition
\[
\varphi(\xx,t)=\int\frac{d^2\kk}{(2\pi)^2}\int_0^\infty\frac{dq}{2\pi}\,
\Bigl[A(\kk,q)\,e^{i\kk\cdot x_\parallel-i\omega t}+\text{c.c.}\Bigr]f_q(z)\,,\qquad
\omega=\sqrt{\kk^2+q^2+m^2}\,.
\]

\emph{(d)} For $M<0$ the decaying solution exists, with $\kappa=-2M=2|M|$:
\[
\varphi=\Re\bigl[A\,e^{i\kk\cdot x_\parallel-i\omega t}\bigr]\,e^{-2|M|z}\,,\qquad\omega^2=\kk^2+m^2-4M^2\,.
\]
These are waves bound to the boundary within a depth $1/2|M|$. If $4M^2<m^2$
they propagate along the surface with the gap $\sqrt{m^2-4M^2}$. If $4M^2>m^2$
then $\omega^2<0$ for $\kk^2<4M^2-m^2$, and the solutions grow as $e^{\Omega t}$
with $\Omega=\sqrt{4M^2-m^2-\kk^2}$. The negative boundary energy
$M\varphi_B^2$ then makes $E$ unbounded below, and the configuration
$\varphi=0$ is unstable.

\noindent\hyperref[hw:1]{$\leftarrow$ Back to Problem 1}

\subsection{Problem 2}\label{sol:2}

Take Dirichlet plates at $x_\perp=0$ and $x_\perp=a$, each of area $A=L^{d-1}$.
As in Section~3, the modes between the plates have
$\omega=\sqrt{\pp_\parallel^2+\pi^2n^2/a^2}$ with $n=1,2,\dots$, and the
regularized zero-point energy per unit area is $\sum_{n\ge1}F(n)$ with
\[
F(\tau)=\frac12\int\frac{d^{d-1}\pp_\parallel}{(2\pi)^{d-1}}\,\sqrt{\pp_\parallel^2+\frac{\pi^2\tau^2}{a^2}}\;\Phi\biggl(\frac{\pp_\parallel^2+\pi^2\tau^2/a^2}{\Lambda^2}\biggr)\,.
\]
The Abel-Plana formula gives
\[
\sum_{n=1}^\infty F(n)=\int_0^\infty F(\tau)\,d\tau-\frac12\,F(0)
+i\int_0^\infty\frac{F(it+0)-F(-it+0)}{e^{2\pi t}-1}\,dt\,.
\]
The first term is proportional to $a$ and is absorbed into the bulk vacuum
energy. The second does not depend on $a$. The last term is finite as
$\Lambda\to\infty$, and it is the Casimir energy. On the two sides of the cut the
square root equals $\pm i\sqrt{\pi^2t^2/a^2-\pp_\parallel^2}$ when
$|\pp_\parallel|<\pi t/a$, and is analytic otherwise. Hence
\[
\frac{E_{\text{Cas}}}A=-\int_0^\infty\frac{dt}{e^{2\pi t}-1}\int_{|\pp|<\pi t/a}\frac{d^{d-1}\pp}{(2\pi)^{d-1}}\,\sqrt{\frac{\pi^2t^2}{a^2}-\pp^2}\,.
\]
The inner integral over the $(d-1)$-ball of radius $R=\pi t/a$ is
\[
\int_{|\pp|<R}\frac{d^{d-1}\pp}{(2\pi)^{d-1}}\,\sqrt{R^2-\pp^2}
=\frac{\pi^{d/2}\,R^d}{2\,(2\pi)^{d-1}\,\Gamma\bigl(\frac d2+1\bigr)}\,,
\]
and $\int_0^\infty t^d\,dt/(e^{2\pi t}-1)=\Gamma(d+1)\,\zeta(d+1)/(2\pi)^{d+1}$.
With the duplication formula for $\Gamma(d+1)$ the result simplifies to
\[
\frac{E_{\text{Cas}}}A=-\frac{\Gamma\bigl(\frac{d+1}2\bigr)\,\zeta(d+1)}{(4\pi)^{\frac{d+1}2}\,a^d}\,,
\qquad
\frac{\text{Force}}A=-\frac{\partial}{\partial a}\,\frac{E_{\text{Cas}}}A
=-\frac{d\;\Gamma\bigl(\frac{d+1}2\bigr)\,\zeta(d+1)}{(4\pi)^{\frac{d+1}2}\,a^{d+1}}\,.
\]
The force is attractive in every dimension. For $d=3$ this gives
$E/A=-\pi^2/1440\,a^3$ and $\text{Force}/A=-\pi^2/480\,a^4$, as in
Section~3. For $d=1$ it gives $E=-\pi/24a$.

\noindent\hyperref[hw:2]{$\leftarrow$ Back to Problem 2}

\subsection{Problem 3}\label{sol:3}

Write $W^{(n)}$ for the connected functions at $J=0$, and expand
\[
W[J]=W[0]+\int_x\bar\phi\,J+\frac12\int W^{(2)}JJ+\frac1{3!}\int W^{(3)}JJJ+\frac1{4!}\int W^{(4)}JJJJ+O(J^5)\,,
\]
with $\int W^{(2)}JJ:=\int d^4x\,d^4y\,W^{(2)}(x,y)\,J(x)J(y)$ and so on. In the
$\varphi^3$ theory $\bar\phi=W^{(1)}\neq0$, which is why the transform is
expanded around $\bar\phi$. Since $\delta\Gamma/\delta\phi=J$ vanishes at
$\phi=\bar\phi$, the expansion of $\hat\Gamma$ starts at second order,
\[
\hat\Gamma[\phi]=\frac12\int\hat\Gamma^{(2)}\phi\phi+\frac1{3!}\int\hat\Gamma^{(3)}\phi\phi\phi+\frac1{4!}\int\hat\Gamma^{(4)}\phi\phi\phi\phi+O(\phi^5)\,.
\]
The Legendre transform pairs the two series through
\begin{align*}
\phi&=\frac{\delta W}{\delta J}-\bar\phi=W^{(2)}J+\frac12\,W^{(3)}JJ+\frac16\,W^{(4)}JJJ+\dots\,,\\
J&=\frac{\delta\hat\Gamma}{\delta\phi}=\hat\Gamma^{(2)}\phi+\frac12\,\hat\Gamma^{(3)}\phi\phi+\frac16\,\hat\Gamma^{(4)}\phi\phi\phi+\dots\,.
\end{align*}
Substitute the second line into the first and collect powers of $\phi$. At first
order $W^{(2)}\hat\Gamma^{(2)}=1$, as before. At second order
\[
W^{(3)}(x_1,x_2,x_3)=-\int_{y_1,y_2,y_3}W^{(2)}(x_1,y_1)\,W^{(2)}(x_2,y_2)\,W^{(2)}(x_3,y_3)\,\hat\Gamma^{(3)}(y_1,y_2,y_3)\,,
\]
and at third order, writing $W$ for $W^{(2)}$,
\begin{multline*}
W^{(4)}(x_1,\dots,x_4)=-\int_{y_1\dots y_4}\prod_{a=1}^4W(x_a,y_a)\;\hat\Gamma^{(4)}(y_1,\dots,y_4)\\
+\biggl[\int_{y_1\dots y_4,u,v}W(x_1,y_1)W(x_2,y_2)\,\hat\Gamma^{(3)}(y_1,y_2,u)\,W(u,v)\,\hat\Gamma^{(3)}(v,y_3,y_4)\,W(y_3,x_3)W(y_4,x_4)\\
+(2\leftrightarrow3)+(2\leftrightarrow4)\biggr]\,.
\end{multline*}
The same results follow from differentiating $W^{(2)}=[\Gamma^{(2)}]^{-1}$ with
respect to $J$: each derivative of $W^{(2)}$ brings $-W^{(2)}W^{(2)}W^{(2)}\Gamma^{(3)}$,
and each derivative of $\Gamma^{(n)}$ brings $\Gamma^{(n+1)}W^{(2)}$.

In diagrams, a hatched blob is $\hat\Gamma^{(n)}$ and a line with an open circle
is the full propagator $W^{(2)}$:
\[
W^{(3)}=-\;\dia[0.4]{\pic{dressed={0.5}{90,210,330}};}\,,\qquad
W^{(4)}=-\;\dia[0.4]{\pic{dressed={0.5}{45,135,225,315}};}
\;+\;\dia[0.4]{\pic{gtree};}\;+\;(\text{two more channels})\,.
\]
At lowest order $\hat\Gamma^{(3)}=\lambda_3$ and $\hat\Gamma^{(4)}$ starts with
the one-loop box, $O(\lambda_3^4)$. So at tree level $W^{(4)}$ is given entirely by
the exchange diagrams.

\noindent\hyperref[hw:3]{$\leftarrow$ Back to Problem 3}

\subsection{Problem 4}\label{sol:4}

All cactus diagrams have momentum-independent self-energy, so write
$\Gamma^{(2)}(p)=p^2+m^2$ with $m^2=m_0^2+\Sigma$. A cactus diagram is a tadpole
whose loop carries any number of smaller cacti, and each of those is built the same
way. Summing the insertions on the loop turns the bare propagator into
$1/(k^2+m_0^2+\Sigma)=1/(k^2+m^2)$, so the sum of all cacti obeys the
self-consistent (Hartree) equation
\[
m^2=m_0^2+\frac\lambda2\int\frac{d^4k}{(2\pi)^4}\,\frac{\Phi(k^2/\Lambda^2)}{k^2+m^2}\,.
\]
Here the cutoff factor multiplies each loop integration, as in
Section~11. With $\Phi=e^{-k^2/\Lambda^2}$ and $u=k^2$,
\[
\int\frac{d^4k}{(2\pi)^4}\,\frac{e^{-k^2/\Lambda^2}}{k^2+m^2}
=\frac1{16\pi^2}\int_0^\infty du\,\frac{u\,e^{-u/\Lambda^2}}{u+m^2}
=\frac1{16\pi^2}\Bigl[\Lambda^2-m^2\,e^{m^2/\Lambda^2}E_1\bigl(m^2/\Lambda^2\bigr)\Bigr]\,,
\]
where $E_1(x)=\int_x^\infty e^{-s}\,ds/s$. The relation, exact in $\lambda$ within
the cactus approximation, is
\[
m_0^2=m^2-\frac{\lambda\,\Lambda^2}{32\pi^2}\Bigl[1-\frac{m^2}{\Lambda^2}\,e^{m^2/\Lambda^2}E_1\Bigl(\frac{m^2}{\Lambda^2}\Bigr)\Bigr]\,.
\]
The function $x\,e^xE_1(x)$ grows monotonically from $0$ to $1$, so for each
$m_0^2$ above the critical value there is exactly one $m^2>0$. For $\Lambda\gg m$,
using $E_1(x)=-\gamma_E-\log x+O(x)$,
\[
m_0^2=m^2-\frac\lambda{32\pi^2}\Bigl[\Lambda^2-m^2\log\frac{\Lambda^2}{m^2}+\gamma_E\,m^2+O\Bigl(\frac{m^4}{\Lambda^2}\log\frac{\Lambda^2}{m^2}\Bigr)\Bigr]\,.
\]
Expanding $m^2$ in powers of $\lambda$ reproduces the first-order result of
Section~11 and then the whole cactus series. The massless theory sits at
$m_0^2=-\lambda\Lambda^2/32\pi^2$.

\noindent\hyperref[hw:4]{$\leftarrow$ Back to Problem 4}

\subsection{Problem 5}\label{sol:5}

Work at $d=4-\epsilon$. The coupling then has dimension $\epsilon$; write it as
$\lambda\mu^\epsilon$ with dimensionless $\lambda=\lambda_{(\mu)}$, so that $\mu$
also sets the scale in the logarithms below. Two massless integrals are needed,
the bubble and the sunset:
\begin{align*}
B(p^2)&=\int\frac{d^dk}{(2\pi)^d}\,\frac1{k^2(k+p)^2}
=\frac{\Gamma(2-\frac d2)\,\Gamma(\frac d2-1)^2}{(4\pi)^{d/2}\,\Gamma(d-2)}\,(p^2)^{\frac d2-2}\\
&=\frac1{(4\pi)^2}\Bigl[\frac2\epsilon+2-\gamma_E+\log4\pi-\log\frac{p^2}{\mu^2}\Bigr]+O(\epsilon)\,,\\
S(p^2)&=\int\frac{d^dk\,d^dq}{(2\pi)^{2d}}\,\frac1{k^2q^2(p-k-q)^2}
=\frac{\Gamma(3-d)\,\Gamma(\frac d2-1)^3}{(4\pi)^d\,\Gamma(\frac{3d}2-3)}\,(p^2)^{d-3}\\
&=-\frac{p^2}{2(4\pi)^4}\Bigl[\frac1\epsilon+\frac{13}4-\gamma_E+\log4\pi-\log\frac{p^2}{\mu^2}\Bigr]+O(\epsilon)\,.
\end{align*}
The sunset follows by doing the $q$ integral first (a bubble with momentum $p-k$)
and then the $k$ integral with the same formula.

\emph{Mass counterterms.} At order $\lambda$ the only diagram in $\Gamma^{(2)}$ is
the tadpole, $\frac\lambda2\int d^dk/k^2=0$. Then \eqref{eq:mu3} gives
$b^{(1)}_{(\mu)}=0$. At order $\lambda^2$ the double tadpole and the tadpole with a
mass counterterm both vanish, and the sunset is proportional to $(p^2)^{d-3}$,
which vanishes at $p=0$. So \eqref{eq:mu3} gives $b^{(2)}_{(\mu)}=0$.

\emph{Two-point vertex.} The sunset enters with two vertices $-\lambda$ and the
symmetry factor $\frac16$:
\[
\Gamma^{(2)}(p^2)=p^2+z^{(2)}_{(\mu)}\lambda^2p^2-\frac{\lambda^2}6\,S(p^2)
=p^2+z^{(2)}_{(\mu)}\lambda^2p^2+\frac{\lambda^2\,p^2}{12(4\pi)^4}\Bigl[\frac1\epsilon+\frac{13}4-\gamma_E+\log4\pi-\log\frac{p^2}{\mu^2}\Bigr]\,.
\]
The condition \eqref{eq:mu1} fixes
\[
z^{(2)}_{(\mu)}=-\frac1{12(4\pi)^4}\Bigl[\frac1\epsilon+\frac94-\gamma_E+\log4\pi\Bigr]\,,
\]
and then
\[
\Gamma^{(2)}(p^2)=p^2+\frac{\lambda^2}{12(4\pi)^4}\,p^2\Bigl(\log\frac{\mu^2}{p^2}+1\Bigr)+O(\lambda^3)\,,
\]
which is \eqref{eq:G2mu}.

\emph{Four-point vertex.} With $p_{ij}^2=(p_i+p_j)^2$,
\[
\Gamma^{(4)}=\lambda+a^{(2)}_{(\mu)}\lambda^2-\frac{\lambda^2}2\Bigl[B(p_{12}^2)+B(p_{13}^2)+B(p_{14}^2)\Bigr]+O(\lambda^3)\,.
\]
At the normalization point $p_i=\mu e_i$ every $p_{ij}^2=\mu^2(e_i+e_j)^2=4\mu^2$, so
\eqref{eq:mu2} fixes
\[
a^{(2)}_{(\mu)}=\frac3{2(4\pi)^2}\Bigl[\frac2\epsilon+2-\gamma_E+\log4\pi-\log4\Bigr]\,,
\]
and
\[
\Gamma^{(4)}(p_1,\dots,p_4)=\lambda+\frac{\lambda^2}{2(4\pi)^2}\left[\log\frac{p_{12}^2}{4\mu^2}+\log\frac{p_{13}^2}{4\mu^2}+\log\frac{p_{14}^2}{4\mu^2}\right]+O(\lambda^3)\,,
\]
which is \eqref{eq:G4mu}. The counterterm $\delta Z$ is $O(\lambda^2)$, so it
enters $\Gamma^{(4)}$ only at order $\lambda^3$.

\noindent\hyperref[hw:5]{$\leftarrow$ Back to Problem 5}

\subsection{Problem 6}\label{sol:6}

\emph{(a)} The measure and the action are invariant under
$\varphi_i\to R_{ij}\varphi_j$ with $R\in O(N)$. So $\tilde W^{(2)}_{ij}$ and
$\tilde W^{(4)}_{ijkl}$ are invariant tensors of $O(N)$. The only invariant tensor
of rank two is $\delta_{ij}$, and those of rank four are
$\delta_{ij}\delta_{kl}$, $\delta_{ik}\delta_{jl}$ and $\delta_{il}\delta_{jk}$.
Bose symmetry of $\tilde W^{(4)}$ under simultaneous exchange of an index and its
momentum relates the three coefficient functions, which gives the stated form
with a single function $\tilde W^{(4)}(p_1,p_2;p_3,p_4)$, symmetric under
$p_1\leftrightarrow p_2$, under $p_3\leftrightarrow p_4$, and under exchange of
the pairs. Since $\tilde W^{(2)}_{ij}$ is diagonal, the proper vertices carry the
same tensors:
\begin{align*}
\Gamma^{(2)}_{ij}(p)&=\delta_{ij}\,\Gamma^{(2)}(p)\,,\qquad\Gamma^{(2)}(p)=1/\tilde W^{(2)}(p)\,,\\
\Gamma^{(4)}_{ijkl}&=\delta_{ij}\delta_{kl}\,\Gamma^{(4)}(p_1,p_2;p_3,p_4)+\delta_{ik}\delta_{jl}\,\Gamma^{(4)}(p_1,p_3;p_2,p_4)+\delta_{il}\delta_{jk}\,\Gamma^{(4)}(p_1,p_4;p_2,p_3)\,,
\end{align*}
where $\tilde W^{(4)}_{ijkl,\,\text{amp}}=-\Gamma^{(4)}_{ijkl}$ as in the
one-component theory.

\emph{(b)} The propagator and the vertex are
\[
\dia[0.5]{\draw[prop] (0,0) -- (2,0); \node[left] at (0,0) {$i$}; \node[right] at (2,0) {$j$}; \draw[arr] (0.7,0.25) -- (1.3,0.25) node[midway,above] {$p$};}
=\frac{\delta_{ij}}{p^2+M^2}\,,\qquad
\dia[0.5]{\pic{xvertex}; \node[left] at (-0.5,0.5) {$i$}; \node[right] at (0.5,0.5) {$j$}; \node[left] at (-0.5,-0.5) {$k$}; \node[right] at (0.5,-0.5) {$l$};}
=-\lambda\,\bigl(\delta_{ij}\delta_{kl}+\delta_{ik}\delta_{jl}+\delta_{il}\delta_{jk}\bigr)\,(2\pi)^d\,\delta^{(d)}(p_1+\dots+p_4)\,.
\]
The vertex is the fourth derivative of $\frac\lambda8(\bm{\varphi}^2)^2$. Each
loop carries $\int d^dk/(2\pi)^d$ and a sum over the internal indices, and the
symmetry factors are those of the one-component theory. For $N=1$ the vertex is
$-3\lambda$, which is $\varphi^4$ theory with $\lambda_{4!}=3\lambda$. At tree
level $\Gamma^{(4)}(p_1,p_2;p_3,p_4)=\lambda$.

\emph{(c)} The scheme is the one of Section~15, applied to the scalar
functions: $d\Gamma^{(2)}/dp^2=1$ at $p^2=\mu^2$, and
$\Gamma^{(4)}(p_1,p_2;p_3,p_4)=\lambda$ at $p_i=\mu e_i$. The renormalization
commutes with $O(N)$, so each tensor structure obeys the equation separately:
\[
\Bigl(\mu\frac\partial{\partial\mu}+\beta(\lambda)\frac\partial{\partial\lambda}-2\gamma(\lambda)\Bigr)\Gamma^{(2)}(p)=0\,,\qquad
\Bigl(\mu\frac\partial{\partial\mu}+\beta(\lambda)\frac\partial{\partial\lambda}-4\gamma(\lambda)\Bigr)\Gamma^{(4)}(p_1,p_2;p_3,p_4)=0\,.
\]

\emph{(d)} Write $V_{ijkl}=\lambda(\delta_{ij}\delta_{kl}+\delta_{ik}\delta_{jl}+\delta_{il}\delta_{jk})$.
The two index sums that appear are
\[
\begin{gathered}
\sum_{m,n}V_{ijmn}V_{mnkl}=\lambda^2\bigl[(N+4)\,\delta_{ij}\delta_{kl}+2\,\delta_{ik}\delta_{jl}+2\,\delta_{il}\delta_{jk}\bigr]\,,\\
\sum_{k,l,m}V_{iklm}V_{jklm}=3(N+2)\,\lambda^2\,\delta_{ij}\,.
\end{gathered}
\]
The one-loop four-point function is $-\frac12\sum_{mn}V_{ijmn}V_{mnkl}\,B(p_{12}^2)$
plus the two crossed channels. Collecting the coefficient of
$\delta_{ij}\delta_{kl}$ and renormalizing as in Problem~5,
\[
\Gamma^{(4)}(p_1,p_2;p_3,p_4)=\lambda+\frac{\lambda^2}{2(4\pi)^2}\left[(N+4)\log\frac{p_{12}^2}{4\mu^2}+2\log\frac{p_{13}^2}{4\mu^2}+2\log\frac{p_{14}^2}{4\mu^2}\right]+O(\lambda^3)\,.
\]
The divergent parts of the three channels add up to
$(N+8)(\delta_{ij}\delta_{kl}+\delta_{ik}\delta_{jl}+\delta_{il}\delta_{jk})$, the
tensor of the bare vertex, so one counterterm $\delta\lambda$ suffices. The sunset
in $\Gamma^{(2)}$ comes with $\frac16\cdot3(N+2)\lambda^2$, so
\[
\Gamma^{(2)}(p)=p^2+\frac{(N+2)\,\lambda^2}{4(4\pi)^4}\,p^2\Bigl(\log\frac{\mu^2}{p^2}+1\Bigr)+O(\lambda^3)\,.
\]
Substituting into the equations of part (c), exactly as in Section~15,
\[
\gamma(\lambda)=\frac{(N+2)\,\lambda^2}{4(4\pi)^4}+O(\lambda^3)\,,\qquad
\beta(\lambda)=\frac{(N+8)\,\lambda^2}{(4\pi)^2}+O(\lambda^3)\,.
\]
For $N=1$ and $\lambda_{4!}=3\lambda$ these become $\gamma=\lambda_{4!}^2/12(4\pi)^4$
and $\beta_{4!}=3\lambda_{4!}^2/(4\pi)^2$, the results \eqref{eq:gamma1} and
\eqref{eq:beta1}.

\noindent\hyperref[hw:6]{$\leftarrow$ Back to Problem 6}

\subsection{Problem 7}\label{sol:7}

\emph{(b)} I do part (b) first. Schwinger's representation and the Gaussian
momentum integral give
\[
D_E(x_E)=\int_0^\infty d\tau\int\frac{d^4p}{(2\pi)^4}\,e^{-\tau(p^2+m^2)+ipx_E}
=\int_0^\infty\frac{d\tau}{(4\pi\tau)^2}\,e^{-m^2\tau-\frac{R^2}{4\tau}}\,.
\]
With $K_\nu(z)=\frac12\bigl(\frac z2\bigr)^\nu\int_0^\infty dt\,t^{-\nu-1}e^{-t-z^2/4t}$
at $\nu=1$ and $t=m^2\tau$,
\[
D_E(x_E)=\frac m{4\pi^2R}\,K_1(mR)\,.
\]

\emph{(a)} From $K_1(z)=\sqrt{\pi/2z}\;e^{-z}\bigl(1+\frac3{8z}+\dots\bigr)$,
\[
D_E(x_E)\simeq\frac1{2(2\pi)^{3/2}}\,\frac{m^{1/2}}{R^{3/2}}\,e^{-mR}\,,\qquad
\Const=\frac1{2(2\pi)^{3/2}}=\frac1{\sqrt{32\pi^3}}\,.
\]
The same constant comes from the saddle point of the $\tau$ integral at
$\tau_*=R/2m$.

\emph{(c)} At small $z$,
$K_1(z)=\frac1z+I_1(z)\log\frac z2+z\,h(z^2)$ with $h$ entire and
$I_1(z)/z$ entire in $z^2$. Write $s^2=x_E^2$. Then
\[
D_E=\frac1{4\pi^2}\left[\frac1{s^2}+\frac12\,\frac{m\,I_1(ms)}s\,\log\frac{m^2s^2}4+E(s^2)\right]
=\frac1{4\pi^2s^2}+\frac{m^2}{16\pi^2}\log\frac{m^2s^2}4+\dots\,,
\]
with $E$ an entire function. The singularity at $x_E^2=0$ is a pole in $s^2$
plus a logarithm multiplied by an entire function. Now set $x_4=it$, so
$s^2=\xx^2-t^2$. As a function of complex $t$, $D(\xx,t)$ is analytic except at
$t=\pm|\xx|$, where it has simple poles and logarithmic branch points. The cuts can
be taken along $t\ge|\xx|$ and $t\le-|\xx|$, where $\xx^2-t^2<0$, as in
Section~4. For real $t>0$ the shifts $t\mp i0$ move $\xx^2-t^2$ to
$\pm i0$. With $1/(u+i0)-1/(u-i0)=-2\pi i\,\delta(u)$, a jump of $2\pi i$ in the
logarithm, and $m\,I_1(ms)/s=m\,J_1(m\sigma)/\sigma$ for
$\sigma=\sqrt{t^2-\xx^2}$,
\[
D(\xx,t-i0)-D(\xx,t+i0)=\operatorname{sign}(t)\left[-\frac i{2\pi}\,\delta(t^2-\xx^2)
+\frac{i\,m}{4\pi}\,\theta(t^2-\xx^2)\,\frac{J_1\bigl(m\sqrt{t^2-\xx^2}\bigr)}{\sqrt{t^2-\xx^2}}\right]\,.
\]
By Section~4 this is
$\bra0[\hat\varphi(\xx,t),\hat\varphi(\mathbf0,0)]\ket0$. It is supported on and
inside the light cone, a delta function on the cone itself plus a Bessel tail
inside, and it vanishes for $|t|<|\xx|$.

\noindent\hyperref[hw:7]{$\leftarrow$ Back to Problem 7}

\subsection{Problem 8}\label{sol:8}

Use the results of Problem~6 with the dimensionless coupling $\lambda\mu^{-\epsilon}$.
At $d=4-\epsilon$ the canonical dimension of the coupling adds $-\epsilon\lambda$
to the flow function:
\[
\beta(\lambda)=-\epsilon\lambda+\frac{(N+8)\,\lambda^2}{16\pi^2}+O(\lambda^3)\,,\qquad
\lambda_*=\frac{16\pi^2\,\epsilon}{N+8}\,.
\]
The anomalous dimension of $\bm{\varphi}^2$ comes from the same one-loop diagram
as in Section~15. The insertion of $\varphi_k\varphi_k$ meets the vertex
through $\sum_k V_{kkij}=(N+2)\lambda\,\delta_{ij}$, so
\[
\gamma_{\varphi^2}(\lambda)=\frac{(N+2)\,\lambda}{16\pi^2}+O(\lambda^2)\,,\qquad
\gamma_{\varphi^2}(\lambda_*)=\frac{N+2}{N+8}\,\epsilon\,.
\]
The anomalous dimension of $\bm{\varphi}$ starts at two loops,
$\gamma(\lambda_*)=\frac{(N+2)\,\epsilon^2}{4(N+8)^2}$, so to order $\epsilon$ it
vanishes. The scaling dimensions at the fixed point are
\[
D_\varphi=\frac{d-2}2+\gamma(\lambda_*)=1-\frac\epsilon2+O(\epsilon^2)\,,\qquad
D_{\varphi^2}=d-2+\gamma_{\varphi^2}(\lambda_*)=2-\frac{6\,\epsilon}{N+8}+O(\epsilon^2)\,.
\]
As in Section~21, $\bm{\varphi}^2$ is the relevant direction, so
$1/\nu=d-D_{\varphi^2}$ and $\alpha=2-d\,\nu$:
\[
\nu=\frac12+\frac{N+2}{4(N+8)}\,\epsilon+O(\epsilon^2)\,,\qquad
\alpha=\frac{4-N}{2(N+8)}\,\epsilon+O(\epsilon^2)\,.
\]
For $N=1$ these reduce to $\nu=\frac12+\frac\epsilon{12}$ and $\alpha=\frac\epsilon6$
of Section~21. For $N\to\infty$ they agree with the large-$N$ values
$\nu=1/(d-2)$ and $\alpha=(d-4)/(d-2)$ to first order in $\epsilon$.

\noindent\hyperref[hw:8]{$\leftarrow$ Back to Problem 8}

\subsection{Problem 9}\label{sol:9}

From Section~22, at small $g$,
\[
\frac{dg}{dl}=\frac{N-2}{2\pi}\,g^2\,,\qquad\gamma_{\mathbf n}(g)=\frac{N-1}{4\pi}\,g\,,\qquad
-\frac d{dl}\,Z_{\mathbf n}^{-\frac12}=\gamma_{\mathbf n}\,Z_{\mathbf n}^{-\frac12}\,,
\]
so $\frac d{dl}\log Z_{\mathbf n}=2\gamma_{\mathbf n}\bigl(g(l)\bigr)$.

\emph{I(a)} The flow gives $1/g(l)=1/g(0)-\frac{N-2}{2\pi}\,l$. Change the variable
from $l$ to $g$, $dl=\frac{2\pi}{N-2}\,dg/g^2$:
\[
\log Z_{\mathbf n}(l)=\frac{N-1}{2\pi}\int_0^lg\,dl'=\frac{N-1}{N-2}\,\log\frac{g(l)}{g(0)}\,,
\]
so
\[
Z_{\mathbf n}(l)=\left(\frac{g(l)}{g(0)}\right)^{\frac{N-1}{N-2}}=\Bigl(1-\frac{N-2}{2\pi}\,g(0)\,l\Bigr)^{-\frac{N-1}{N-2}}\,.
\]

\emph{I(b)} Take the scaling relation of Section~22 with $|x_1-x_2|=r$ and
$L=r$, so $l=\log r<0$:
\[
\bigl\langle\mathbf n(x)\,\mathbf n(0)\bigr\rangle_{g(0)}\Big|_{|x|=r}
=Z_{\mathbf n}^{-1}(l)\,\bigl\langle\mathbf n(x)\,\mathbf n(0)\bigr\rangle_{g(l)}\Big|_{|x|=1}\,.
\]
As $r\to0$ the coupling $g(l)\simeq2\pi/\bigl((N-2)\log\frac1r\bigr)$ goes to zero,
and the correlator on the right tends to a constant. Therefore
\[
\bigl\langle\mathbf n(x)\,\mathbf n(0)\bigr\rangle\simeq\Const\,\Bigl(\frac{g(0)}{g(l)}\Bigr)^{\frac{N-1}{N-2}}
\sim\left[\log\frac1{x^2}\right]^{\frac{N-1}{N-2}}\qquad(|x|\to0)\,.
\]

\emph{I(c)} At $N\to\infty$ the exponent tends to $1$. In the exact large-$N$
solution (PS, \S13.3) each component of $\mathbf n$ is a free field of mass
$m\propto\Lambda e^{-2\pi/Ng}$, with propagator $1/(p^2+m^2)$. In two dimensions
its Fourier transform is $\frac1{2\pi}K_0(m|x|)\simeq\frac1{4\pi}\log\frac1{m^2x^2}$
at small $|x|$, a single power of the logarithm, as the RG result says.

\emph{II} At $d=2+\epsilon$,
\[
\frac{dg}{dl}=-\epsilon g+\frac{N-2}{2\pi}\,g^2=\frac{N-2}{2\pi}\,g\,(g-g_*)\,,\qquad
\frac1{g(l)}=\frac1{g_*}+\Bigl(\frac1{g(0)}-\frac1{g_*}\Bigr)e^{\epsilon l}\,,\qquad g_*=\frac{2\pi\epsilon}{N-2}\,,
\]
and the flow equation for $Z_{\mathbf n}$ is
\[
\frac d{dl}\,Z_{\mathbf n}(l)=\frac{N-1}{2\pi}\,g(l)\,Z_{\mathbf n}(l)\,,\qquad
Z_{\mathbf n}(l)=\left(\frac{g(l)-g_*}{g(0)-g_*}\right)^{\frac{N-1}{N-2}}\,.
\]
At $\epsilon\to0$ this goes back to part I(a). Going to short distances,
$l\to-\infty$, the coupling flows to the ultraviolet fixed point $g_*$ from either
side, with $g(l)-g_*\propto e^{\epsilon l}$. Then
$Z_{\mathbf n}(l)\propto e^{2\gamma_*l}$ with
$\gamma_*=\gamma_{\mathbf n}(g_*)=\frac{(N-1)\,\epsilon}{2(N-2)}$, and the same
argument as in I(b) gives a power law,
\[
\bigl\langle\mathbf n(x)\,\mathbf n(0)\bigr\rangle\sim|x|^{-2\gamma_*}=|x|^{-\frac{N-1}{N-2}\,\epsilon}\qquad(|x|\to0)\,.
\]
In the usual notation $2D_{\mathbf n}=d-2+\eta$ this is $\eta=\epsilon/(N-2)$, the
known result for the sigma model at $2+\epsilon$ dimensions.

\noindent\hyperref[hw:9]{$\leftarrow$ Back to Problem 9}
